\documentclass{../../../unipd-notes}

\unipdsetup{
  course = {Fondamenti di Ingegneria Informatica},
  short-course = {Fondamenti di Ingegneria},
  document-type = {Appunti delle lezioni}
}

\begin{document}
\chapter{Ambienti didattici}

\begin{definition}[Numero pari]\label{def:even}
Un intero $n$ è pari se esiste $k\in\Z$ tale che $n=2k$.
\end{definition}

\begin{theorem}[Somma di numeri pari]\label{thm:sum}
La somma di due numeri pari è pari.
\end{theorem}

\begin{proof}
Se $a=2h$ e $b=2k$, allora $a+b=2(h+k)$.
\end{proof}

\begin{proposition}\label{prop:product}
Il prodotto di un numero pari per un intero è pari.
\end{proposition}

\begin{lemma}\label{lem:multiple}
Ogni intero della forma $2k$ è divisibile per due.
\end{lemma}

\begin{corollary}\label{cor:double}
Il doppio di ogni intero è pari.
\end{corollary}

\begin{example}\label{ex:even}
La somma di $8$ e $14$ è $22$.
\end{example}

\begin{remark}\label{rem:numbering}
Ogni ambiente numerato può ricevere un titolo e un'etichetta.
\end{remark}

\begin{important}[Idea da ricordare]\label{imp:idea}
Le informazioni centrali possono essere evidenziate con moderazione.
\end{important}

\begin{warning}[Ipotesi necessarie]\label{warn:hypothesis}
Prima di applicare un risultato occorre verificarne le ipotesi.
\end{warning}

\begin{exercise}\label{exe:sum}
Dimostrare che la somma di quattro numeri pari è pari.
\end{exercise}

\begin{solution}\label{sol:sum}
Si applica due volte il \cref{thm:sum}.
\end{solution}

I riferimenti configurati comprendono \cref{def:even,thm:sum,prop:product,lem:multiple,cor:double,ex:even,rem:numbering,imp:idea,warn:hypothesis,exe:sum,sol:sum}.
\end{document}
